\section{模型评价与推广}

\subsection{模型优点}

\begin{enumerate}
  \item \textbf{方法科学}：问题一针对纵向面板数据采用组内/组间分解与混合效应模型，正确分离了孕周效应与 BMI 效应，避免了横截面混杂导致的错误结论。
  \item \textbf{估计稳健}：达标时间采用``共享斜率 + 个体水平''的稳健估计，克服了少数采样点直接回归的数值不稳定。
  \item \textbf{风险建模直观}：以``组内 90\% 可靠达标''的最早时点为最佳时点，兼顾检测可靠性与早检要求，可解释性强。
  \item \textbf{判定方法可落地}：问题四用随机森林综合多维特征，显著优于单一 Z 值阈值法，且给出可解释的特征重要性。
\end{enumerate}

\subsection{模型不足}

\begin{enumerate}
  \item 极端 BMI 组（$[20,28)$ 与 $\ge 40$）样本量少，最佳时点估计的稳定性较低。
  \item 达标时间估计依赖``共享斜率''假设，若个体间孕周增长率差异较大，会产生一定偏差。
  \item 问题四异常样本量较少（67 例）且均为出生后健康（存在假阳性），判定模型灵敏度仍有限，需更大规模、含真阳性样本的数据进一步验证。
  \item 检测误差 $\sigma_e$ 取混合模型残差标准差作为近似，未区分测量误差与生物变异。
\end{enumerate}

\subsection{模型推广}

\begin{enumerate}
  \item 本文的``面板分解 + 分组风险最小化''框架可推广到其他依赖纵向生物标志物的检测时点优化问题，如其他产前筛查指标、肿瘤标志物监测等。
  \item 女胎异常判定的多维分类方法可推广到更广泛的医学诊断场景，通过引入更多质量指标与更大的临床样本提升灵敏度。
  \item 分组与最佳时点方法可进一步扩展为个性化方案：用孕妇个体特征直接预测其达标时间，给出个体化 NIPT 时点建议。
\end{enumerate}
